%% !TEX root = ../manual_gaolv5.tex
% The conclusion (unnumbered).
\schapter{Conclusion}\label{chap:conclusion}
\thispagestyle{plain}

This manual described GAOL v5, the continuation of the GAOL of Fr\'ed\'eric
Goualard, as a user's guide and as a reference. We recalled the notions it
rests on, from the roundings of floating-point numbers to the accuracy modes
of IEEE 1788-2015; we described its installation with the three builds, the
flags the code using it is compiled with, the two namespaces under which its
operations are named, and the rounding direction and the floating-point
environment that GAOL v5 takes charge of; we gave each operation of the class
\code{interval} with its examples, the names of the standard, and the
tightness of every operation; we explained how the intervals are computed
and measured, how GAOL is tested, and what a program of GAOL~4 has to change.

\medskip
Three facts sum up what GAOL v5 established. First, the bounds it computes
are the same on every architecture and with every compiler it supports:
every elementary function that CORE-MATH provides is bounded with it, correctly
rounded in the direction in effect, and the integer powers, the roots and
\code{compoundm1} are the tightest bounds or within a few doubles of them; the
tests compare the bounds with their exact
values, computed apart, on every platform of the continuous integration.
Second, the operations follow IEEE 1788-2015 for the empty set, the infinite
bounds, the domains of the functions, the comparisons and the literals; under
the names of the standard, GAOL v5 gives the result of the standard in each of
the 288 special cases of the comparison with the other libraries where the standard
defines one (Chapter~\ref{chap:computing}), and under its
own names in all but those of its hybrid \code{pow}, kept for the programs of
GAOL~4. Third, the rounding direction is no longer left to the user: each
operation checks it and sets it upward when it is not, but the operations of
the AVX-512 path, whose instructions carry the direction, and GAOL can be built
to restore the direction of the program after each operation, at a cost that
the AVX-512 path, where the machine has the instructions, brings back to a
few tens of percent.

\medskip
What GAOL v5 does not provide is stated where it would be expected. GAOL has
no decorations: only bare intervals are provided, and a program that needs
to know that a box lies outside the domain of a function tests the empty
set itself, as Chapter~\ref{chap:overview} explains. The slope
functions, the reduction operations and the exact reduction operations of the
standard are not provided either. The intervals of floats are unfinished, and no build
compiles them. The \code{pow} of the namespace \code{gaol} is a hybrid of the
two powers of the standard, which keeps the results of GAOL~4 for the
programs written for it and for IBEX and Codac; a program that wants the
results of the standard opens \code{gaol\_ieee1788}.

\medskip
What is left to measure concerns the processors on which GAOL v5 could not be timed.
Of the instruction sets read for Chapter~\ref{chap:computing}, the scalar
instructions of RISC-V are the only ones, beside the AVX-512 path of x86,
that carry their rounding direction; a path written with them would need
neither a check nor a restore of the direction, and waits for a board to be
measured on. The vector units of POWER and s390x hold two doubles per
register as NEON does, and wait for hardware as well. Until then, the FPU
intervals compute on these processors, with the same bounds, as the
continuous integration checks under emulation.
